% =====================================================================
%  13  水素分離 — PSA（2段：上=概念図/吸着等温線、下=設計変数表/方式選定）
%  source_memo: 08章＋psa_system.py＋top1_trial201.txt[PSA]。活性炭・多成分Langmuir(Choi2003)
%   ＋LDF・1D PDE。H2非吸着で全量製品、CH4/C2H4/C2H6を吸着。破過基準0.1%。
%   最適化変数: 塔径D4.49m・層高L24.8m・脱着目標β0.254。導出 t_abs742/t_des113s・2基。
%   H2回収~84%・H2純度>99.9%・CH4捕捉100%。入1486(H2 60%)→H2製品746＋オフガス740 kmol/h。
%  方針：図で一目＋口頭。H2=緑、吸着成分=暖色。丸括弧禁止。余白を作らず最下部まで。
% =====================================================================
\definecolor{SepGreen}{HTML}{2E9E5B}
\definecolor{DipeOrange}{HTML}{E8821E}
\begin{frame}[t]

  \vspace{0.1em}

  % ===== 上段：PSA概念図（左）／吸着等温線（右）=================
  \begin{columns}[T,onlytextwidth]
    \begin{column}{0.46\textwidth}
      \raggedright
      {\footnotesize\bfseries\color{HeaderBlue} PSA　圧力スイング吸着}\par
      \vspace{0.3em}
      \centering
      \resizebox{\linewidth}{!}{%
      \begin{tikzpicture}[
        x=1cm, y=1cm, >=Stealth,
        bed/.style={draw, line width=1pt, minimum width=1.4cm, minimum height=2.1cm,
                    align=center, font=\Large\bfseries, rounded corners=2pt},
        io/.style={align=center, font=\Large},
        a/.style={-{Stealth[length=2.8mm]}, line width=1.4pt},
      ]
        \node[bed, draw=HeaderBlue, fill=HiBlue!45, text=HeaderBlue] (A) at (0,0) {吸着\\高圧};
        \node[bed, draw=black!45, fill=black!8, text=black!55] (B) at (2.9,0) {脱着\\低圧};
        % 供給（左→A）
        \node[io, anchor=east] (feed) at (-1.55,0) {供給 $\mathrm{H_2}$ リッチ\\{\normalsize $1486$ kmol/h ・$60$\%}};
        \draw[a, HeaderBlue] (feed.east) -- (A.west);
        % H2 製品（A→上）
        \node[io, anchor=south, text=SepGreen!45!black] (h2) at (0,1.5) {$\mathrm{H_2}$ 製品 $>99.9$\%\\{\normalsize $746$ kmol/h}};
        \draw[a, SepGreen!70!black] (A.north) -- (h2.south);
        % オフガス（B→下）
        \node[io, anchor=north, text=DipeOrange!75!black] (off) at (2.9,-1.45) {オフガス → 燃料\\{\normalsize $740$ kmol/h}};
        \draw[a, DipeOrange!80!black] (B.south) -- (off.north);
        % スイング
        \draw[{Stealth[length=2mm]}-{Stealth[length=2mm]}, line width=0.9pt, black!55]
              (A.east) -- node[above, font=\normalsize, text=black!60]{スイング}(B.west);
      \end{tikzpicture}}
    \end{column}

    \begin{column}{0.54\textwidth}
      \centering
      \includegraphics[width=\linewidth]{psa_isotherm}\par
      \vspace{-0.82em}
      {\scriptsize 活性炭の吸着等温線　$25\,^\circ$C}\par
      \vspace{0.12em}
      {\tiny\color{black!55} 出典：B.-U.\ Choi \textit{et al.}, J.\ Chem.\ Eng.\ Data, \textbf{48}(3), 603--607 (2003)}
    \end{column}
  \end{columns}

  \vspace{0.7em}

  % ===== 下段：設計変数表（左）／方式の選定 枠囲み（右）=========
  \begin{columns}[T,onlytextwidth]
    \begin{column}{0.37\textwidth}
      {\scriptsize\bfseries\color{HeaderBlue} 設計変数と諸元}\par
      \vspace{0.12em}
      {\scriptsize\setlength{\tabcolsep}{5pt}\renewcommand{\arraystretch}{1.02}%
       \begin{tabular}{@{}l r@{}}
         \toprule
         塔径 $D$ & $4.49$ m \\
         層高 $L$ & $24.8$ m \\
         脱着目標 $\beta$ & $0.254$ \\
         \midrule
         吸着材 & 活性炭 \\
         塔数 & $2$ 基 \\
         吸着／脱着時間 & $742$／$113$ s \\
         \midrule
         $\mathrm{H_2}$ 純度 & $99.9$\% 超 \\
         $\mathrm{H_2}$ 回収率 & 約 $84$\% \\
         $\mathrm{CH_4}$ 捕捉 & $100$\% \\
         \bottomrule
       \end{tabular}}\par
      \vspace{0.12em}
      {\tiny\color{black!55} 上3行が最適化変数}
    \end{column}

    \begin{column}{0.63\textwidth}
      {\setlength{\fboxsep}{6pt}\setlength{\fboxrule}{1.2pt}%
       \fcolorbox{HeaderBlue!70!black}{HiBlue!12}{%
        \begin{minipage}[t][4.13cm]{0.93\linewidth}
         {\scriptsize\bfseries\color{HeaderBlue} 方式の選定　PSA 圧力スイング吸着}
         \vfill
         {\scriptsize\setlength{\tabcolsep}{4pt}\renewcommand{\arraystretch}{1.1}%
          \begin{tabular}{@{}l c l@{}}
            \toprule
            方式 & $\mathrm{H_2}$ 純度 & 課題 \\
            \midrule
            深冷分離 & 高 & $-170\,^\circ$C 級・高コスト \\
            膜 & 中 & 高純度に多段が必要 \\
            {\bfseries\color{HeaderBlue} 採用 PSA} & {\bfseries $>99.9$\%} & 工業実績・連続 2 基 \\
            \bottomrule
          \end{tabular}}
         \vfill
         {\scriptsize\raggedright 吸着材＝\textbf{活性炭}：$\mathrm{H_2}$ 非吸着・$\mathrm{CH_4}/\mathrm{C_2}$ を選択吸着\par}
         \vfill
         {\scriptsize\raggedright 価値化：$\mathrm{H_2}$ は製品販売・\textbf{オフガスは予熱炉の燃料}\par}
        \end{minipage}}}
    \end{column}
  \end{columns}

\end{frame}
